\begin{figure}[t]
\begin{minted}[fontsize=\small,frame=single]{rust}
// Silq teleportation, corrected
def teleport(msg: B, target: B) : B {
    alice := 0:B;  // allocate |0>

    // create entanglement (not qfree)
    alice := H(alice);
    (alice, target) := CNOT(alice, target);

    // Alice's Bell measurement
    (msg, alice) := CNOT(msg, alice);
    msg := H(msg);

    // classical results
    m1 := measure(msg);     // m1: !B
    m2 := measure(alice);   // m2: !B

    // classical control using measured bits
    if m2 == 1 { target := X(target); }
    if m1 == 1 { target := Z(target); }

    return target;
    // msg and alice were consumed by measure, so nothing to clean up
}
\end{minted}
\caption{Quantum teleportation in Silq. The type \texttt{!B} denotes a classical bit that can be copied and used in classical control. Variables like \texttt{alice} and \texttt{msg} are automatically uncomputed when they go out of scope because their transformations are qfree. The \texttt{measure} operation converts quantum \texttt{B} to classical \texttt{!B}.}
\label{fig:silq-teleport}
\end{figure}
